
\documentclass[11pt]{article}
\usepackage[ngerman]{babel}
\usepackage[a4paper,margin=0.9in]{geometry}
\usepackage[T1]{fontenc}
\usepackage[utf8]{inputenc}
\usepackage{lmodern}
\usepackage{amsmath,amssymb,mathtools,mathrsfs}
\usepackage{microtype}
\usepackage{fancyhdr}
\usepackage{array,longtable,enumitem,multicol}
\usepackage{tikz}
\usetikzlibrary{patterns}
\setlength{\parindent}{1.5em}
\setlength{\parskip}{0.18em}
\emergencystretch=6em
\setlength{\headheight}{14pt}
\pagestyle{fancy}
\fancyhf{}
\cfoot{\thepage}
\providecommand{\sect}[2]{\section*{\S\,#1. #2}}
\providecommand{\dd}{\,d}
\providecommand{\ii}{\mathrm{i}}
\providecommand{\disc}{\operatorname{disc}}
\providecommand{\rank}{\operatorname{rank}}
\providecommand{\Bez}{\mathcal B}

\lhead{Weber Vol. I \S\S 76--81}
\rhead{Deutsche Quelle}
\begin{document}

\begin{center}
{\Large ZWEITES BUCH.}\\[1.6em]
{\LARGE DIE WURZELN.}
\end{center}
\clearpage

\begin{center}
{\large Siebenter Abschnitt.}\\[0.6em]
{\Large Realit\"at der Wurzeln.}
\end{center}

\sect{76}{Allgemeines \"uber Realit\"at von Gleichungswurzeln und \"uber die Discriminante}

In diesem Abschnitt werden wir uns mit der Frage besch\"aftigen, wie viele Wurzeln einer algebraischen Gleichung reell sind. Die Coefficienten der Gleichung werden dabei als reelle Zahlen vorausgesetzt. Wir wollen solche Gleichungen kurz reelle Gleichungen nennen und beginnen mit einigen allgemeinen Betrachtungen.

Die reelle Gleichung $n$ten Grades
\[
        f(x)=0
\]
hat, wie wir im dritten Abschnitt gesehen haben, $n$ Wurzeln, die entweder reell oder imagin\"ar, d. h. von der Form $\xi+\ii\eta$ sind, mit reellen $\xi,\eta$. Die Zahl von $n$ Wurzeln ergiebt sich aber nur dann allgemein, wenn wir unter Umst\"anden eine Wurzel mehrfach z\"ahlen, n\"amlich $(m+1)$fach, wenn mit $f(x)$ zugleich die $m$ ersten Derivirten von $f(x)$ verschwinden. Da nun ein complexer Ausdruck von der Form $X+\ii Y$ nur dann gleich Null ist, wenn die beiden reellen Bestandtheile $X$ und $Y$ einzeln verschwinden, wenn also mit $X+\ii Y$ zugleich $X-\ii Y$ verschwindet, und da ferner bei reellen Coefficienten, wenn $f(\xi+\ii\eta)=X+\ii Y$ ist, sich $f(\xi-\ii\eta)=X-\ii Y$ ergiebt, so folgt aus $f(\xi+\ii\eta)=0$, dass zugleich $f(\xi-\ii\eta)=0$ sein muss, dass also zu jeder imagin\"aren Wurzel $\xi+\ii\eta$ eine zweite davon verschiedene imagin\"are $\xi-\ii\eta$, d. h. die conjugirt imagin\"are Wurzel geh\"ort. Da mit den Derivirten $f'(\xi+\ii\eta), f''(\xi+\ii\eta),\ldots$ zugleich die conjugirten Gr\"ossen $f'(\xi-\ii\eta), f''(\xi-\ii\eta),\ldots$ verschwinden, so folgt, dass conjugirt imagin\"are Wurzeln denselben Grad der Vielfachheit haben.

Daraus folgt, dass die imagin\"aren Wurzeln immer in gerader Zahl vorkommen, und dass eine Gleichung ungeraden Grades immer mindestens eine reelle Wurzel haben muss.

Unser n\"achstes Ziel wird das sein, die Zahl der reellen Wurzeln, ohne die Gleichung aufzul\"osen, direct aus den Werthen gewisser rationaler Functionen der Coefficienten der Gleichung zu bestimmen.

Eine sehr wichtige Rolle spielt hierbei die Discriminante, auf deren Bedeutung f\"ur unsere Frage wir zun\"achst eingehen m\"ussen.

Wir haben in \S\,46 die Discriminante erkl\"art als das Product aus den Quadraten der s\"ammtlichen Wurzeldifferenzen
\[
 (x_1-x_2)^2(x_1-x_3)^2\cdots (x_{n-1}-x_n)^2,
\]
noch multiplicirt mit $a_0^{2n-2}$, wenn $a_0$ der Coefficient der h\"ochsten Potenz der Unbekannten in der gegebenen Gleichung ist, und wir haben dort gezeigt, wie die Discriminante als rationale Function der Coefficienten berechnet werden kann. Der Factor $a_0^{2n-2}$ ist immer positiv und k\"onnte hier ohne wesentliche Beschr\"ankung der Allgemeinheit auch gleich $1$ vorausgesetzt werden. Wir wollen, wie schon fr\"uher, f\"ur die Discriminante das Zeichen $D$ gebrauchen.

Die Discriminante verschwindet dann und nur dann, wenn unter den Wurzeln zwei gleiche vorkommen.

Nehmen wir an, dass $f(x)$ durch Absonderung des gr\"ossten gemeinschaftlichen Theilers von $f(x)$ und $f'(x)$, was durch rationale Rechnung geschieht, von mehrfachen Factoren befreit sei, so wird also $D$ nicht verschwinden.

Sind $x_1$ und $x_2$ reell, so ist $(x_1-x_2)^2$ positiv. Ist $x_1$ reell und $x_2$ imagin\"ar, so giebt es eine zu $x_2$ conjugirte Wurzel $x'_2$, und das Product
\[
        (x_1-x_2)(x_1-x'_2)
\]
ist als Product zweier conjugirt imagin\"arer Gr\"ossen positiv, also auch sein Quadrat.

Sind $x_1$ und $x_2$ beide imagin\"ar, aber nicht conjugirt, so ist das Product
\[
        (x_1-x_2)(x'_1-x'_2)
\]
und sein Quadrat positiv.

Sind aber endlich $x_1$ und $x_2$ conjugirt imagin\"ar, so ist ihre Differenz rein imagin\"ar und deren Quadrat negativ.

Es kommen also in dem Product, durch das wir die Discriminante erkl\"art haben, genau so viel negative Factoren vor, als es Paare conjugirt imagin\"arer Wurzeln giebt, und wir schliessen daraus auf den wichtigen Fundamentalsatz:

Ist $f(x)=0$ eine reelle Gleichung ohne mehrfache Wurzeln, so ist die Discriminante positiv oder negativ, je nachdem die Anzahl der Paare conjugirt imagin\"arer Wurzeln gerade oder ungerade ist.

Der Fall, wo nur reelle Wurzeln vorhanden sind, geh\"ort zu denen, wo die Discriminante positiv ist.

F\"ur die F\"alle der quadratischen und cubischen Gleichungen, in denen nur ein Paar imagin\"arer Wurzeln auftreten kann, ist durch diesen Hauptsatz bereits die Unterscheidung der verschiedenen F\"alle, die in Bezug auf die Wurzelrealit\"at m\"oglich sind, die Determination, vollendet.

\sect{77}{Discussion der quadratischen und cubischen Gleichung}

F\"ur die quadratische Gleichung
\[
        a_0x^2+a_1x+a_2=0
\]
haben wir (\S\,46)
\[
        D=a_1^2-4a_0a_2
\begin{cases}
>0, & \text{reelle Wurzeln,}\cr
<0, & \text{imagin\"are Wurzeln.}
\end{cases}
\]
F\"ur die cubische Gleichung
\[
        a_0x^3+a_1x^2+a_2x+a_3=0
\]
ist
\[
 D=a_1^2a_2^2+18a_0a_1a_2a_3-4a_0a_2^3-4a_1^3a_3-27a_0^2a_3^2,
\]
und
\[
\begin{array}{ll}
D>0, & \text{drei reelle Wurzeln,}\cr
D<0, & \text{eine reelle, zwei imagin\"are Wurzeln.}
\end{array}
\]

Auch der Fall $D=0$ giebt hier zu keinen weiteren Unterscheidungen Anlass; denn in diesem Falle sind die beiden Wurzeln der quadratischen Gleichung einander gleich und reell, von den drei Wurzeln der cubischen Gleichung zwei einander gleich und alle drei reell; denn eine imagin\"are Wurzel kann hier nicht doppelt vorkommen, weil sonst auch die conjugirte doppelt vorkommen w\"urde, und ebenso wenig kann die einzelne Wurzel imagin\"ar sein, weil sonst eine zweite vorhanden sein m\"usste.

Es kann sich bei der cubischen Gleichung nur noch darum handeln, die Bedingung daf\"ur aufzusuchen, dass alle drei Wurzeln einander gleich sind.

In diesem Falle muss die linke Seite der cubischen Gleichung ein vollst\"andiger Cubus sein, also
\[
        a_0x^3+a_1x^2+a_2x+a_3=a_0(x-\alpha)^3.
\]
Die Vergleichung beider Seiten dieser identischen Gleichung ergiebt
\[
        a_1=-3a_0\alpha,\qquad a_2=3a_0\alpha^2,\qquad a_3=-a_0\alpha^3,
\]
woraus man durch Elimination von $\alpha$ erh\"alt
\begin{equation}
        a_1^2-3a_0a_2=0,\qquad a_1a_2-9a_0a_3=0.
\tag{1}
\end{equation}
Daraus ergiebt sich als Folge, indem man die erste dieser Gleichungen mit $a_2$, die zweite mit $a_1$ multiplicirt und subtrahirt,
\[
        a_2^2-3a_1a_3=0;
\]
und wenn diese Bedingungen erf\"ullt sind, so folgt daraus umgekehrt
\[
        f(x)=a_0\left(x+{a_1\over 3a_0}\right)^3,
\]
d. h. die Gleichheit aller drei Wurzeln.

Bei der cubischen Gleichung
\[
        x^3+a_1x^2+a_2x+a_3=0
\]
k\"onnen wir ausser \"uber die Realit\"at auch noch \"uber die Vorzeichen der Wurzeln vollst\"andig entscheiden.

Bezeichnen wir n\"amlich mit $\alpha,\beta,\gamma$ die drei Wurzeln, so ist
\begin{equation}
        a_1=-(\alpha+\beta+\gamma),\qquad
        a_2=\alpha\beta+\alpha\gamma+\beta\gamma,\qquad
        a_3=-\alpha\beta\gamma.
\tag{2}
\end{equation}
Ist nun $D<0$, also nur eine Wurzel, etwa $\alpha$, reell, so ist $\beta\gamma$ positiv und $\alpha$ wird negativ oder positiv sein, je nachdem $a_3$ positiv oder negativ ist. Ist aber $D$ positiv, also sind alle drei Wurzeln reell, so ist, wenn $\alpha,\beta,\gamma$ positiv sind, nach (2)
\begin{equation}
        a_1<0,
        \qquad a_2>0,
        \qquad a_3<0;
\tag{3}
\end{equation}
diese Bedingungen sind nothwendig daf\"ur, dass alle drei Wurzeln positiv sind. Sie sind aber auch hinreichend; denn wenn eine oder drei Wurzeln negativ sind, so ist $a_3>0$. Sind aber zwei Wurzeln, etwa $\beta,\gamma$, negativ, so ist entweder
\[
        a_1\geq 0 \qquad\text{oder}\qquad \alpha>-(\beta+\gamma).
\]
In letzterem Falle aber folgt
\[
        a_2<\beta\gamma-(\beta+\gamma)^2=-\beta^2-\beta\gamma-\gamma^2,
\]
also $a_2$ negativ.

Ist endlich eine Wurzel gleich $0$, so muss nothwendig $a_3$ verschwinden.

Indem man $x$ durch $-x$ ersetzt, schliesst man, dass f\"ur drei negative Wurzeln die nothwendige und hinreichende Bedingung die ist
\begin{equation}
        a_1>0,
        \qquad a_2>0,
        \qquad a_3>0.
\tag{4}
\end{equation}

Wir erhalten daher unter Voraussetzung einer positiven Discriminante folgende Tabelle:
\[
\begin{array}{ccc|c}
 a_1 & a_2 & a_3 & \text{Zahl der positiven Wurzeln}\cr\hline
 - & + & - & 3\cr\hline
 - & - & - & 1\cr
 + & + & - & 1\cr
 + & - & - & 1\cr\hline
 + & + & + & 0\cr\hline
 - & - & + & 2\cr
 - & + & + & 2\cr
 + & - & + & 2
\end{array}
\]
wo in der letzten Columne die Zahl der positiven Wurzeln angegeben ist. Wir k\"onnen das Resultat dieser Betrachtung so aussprechen:

Bei positiver Discriminante ist die Anzahl der positiven Wurzeln der cubischen Gleichung gleich der Anzahl der Zeichenwechsel in der Reihe
\[
        1,a_1,a_2,a_3,
\]
wenn wir unter einem Zeichenwechsel die Aufeinanderfolge einer positiven und einer negativen oder einer negativen und einer positiven Gr\"osse verstehen.

Wenn eine der beiden Gr\"ossen $a_1,a_2$ verschwindet, so haben wir, da dann nicht alle Wurzeln von gleichem Zeichen sein k\"onnen, eine oder zwei positive Wurzeln.

Wenn $a_3=0$ ist, so reducirt sich die Discriminante auf
\[
        a_2^2(a_1^2-4a_2),
\]
und wenn diese positiv ist, so hat die cubische Gleichung eine verschwindende und noch zwei andere reelle Wurzeln. Diese sind positiv, wenn
\[
        a_1<0,\qquad a_2>0,
\]
negativ, wenn
\[
        a_1>0,\qquad a_2>0,
\]
und es ist eine von ihnen positiv und eine negativ, wenn
\[
        a_1\gtreqless 0,\qquad a_2<0
\]
ist.

\sect{78}{Discussion der biquadratischen Gleichung}

Bei den Gleichungen vierten und f\"unften Grades existiren entweder keine oder zwei oder vier imagin\"are Wurzeln. Ist die Discriminante negativ, so hat man zwei imagin\"are Wurzeln. Bei positiver Discriminante k\"onnen entweder vier oder keine imagin\"aren Wurzeln vorhanden sein. Diese beiden F\"alle zu unterscheiden, wird weiterhin unsere Aufgabe sein. Wir wenden uns aber zun\"achst zu einer elementaren Betrachtung der biquadratischen Gleichung, die wir der Einfachheit halber in der verk\"urzten Form
\begin{equation}
        x^4+ax^2+bx+c=0
\tag{1}
\end{equation}
annehmen wollen, von der wir leicht zur allgemeinen Form zur\"uckkehren k\"onnen.

Bezeichnen wir die Wurzeln mit $\alpha,\beta,\gamma,\delta$ und setzen, wie in \S\,47,
\begin{equation}
\begin{array}{ll}
\alpha-\beta=v+w, & \gamma-\delta=v-w,\cr
\alpha-\gamma=w+u, & \delta-\beta=w-u,\cr
\alpha-\delta=u+v, & \beta-\gamma=u-v,
\end{array}
\tag{2}
\end{equation}
so sind $u^2,v^2,w^2$ die Wurzeln der cubischen Resolvente
\begin{equation}
        y^3+2ay^2+(a^2-4c)y-b^2=0,
\tag{3}
\end{equation}
und die Discriminante $D$ dieser cubischen Gleichung, die durch
\begin{equation}
        27D=4(a^2+12c)^3-(2a^3-72ac+27b^2)^2
\tag{4}
\end{equation}
bestimmt ist, ist zugleich die Discriminante der biquadratischen Gleichung (1), und die Vorzeichen der Gr\"ossen $u,v,w$ sind durch die Bedingung
\begin{equation}
        uvw=-b
\tag{5}
\end{equation}
beschr\"ankt [\S\,37, (3)].

Wenn die Discriminante negativ ist, so hat die Gleichung (3) ebenso wie (1) zwei conjugirt imagin\"are Wurzeln.

Ist $D$ positiv und alle vier Wurzeln $\alpha,\beta,\gamma,\delta$ reell, so werden auch $u,v,w$ reell und also ihre Quadrate, d. h. die Wurzeln von (3), positiv.

Sind alle vier Wurzeln imagin\"ar, etwa $\alpha$ mit $\beta$ und $\gamma$ mit $\delta$ conjugirt, so folgt aus (2), dass $v$ und $w$ rein imagin\"ar, $u$ reell ist; also hat in diesem Falle die Gleichung (3) eine positive und zwei negative Wurzeln.

In beiden F\"allen kann aber auch, wenn $b$ verschwindet, eine der Gr\"ossen $u,v,w$ gleich Null sein.

Mit R\"ucksicht auf die Ergebnisse des letzten Paragraphen \"uber die cubische Gleichung kommen wir also zu folgendem Resultat:

Die nothwendige und hinreichende Bedingung f\"ur die Existenz von vier verschiedenen reellen Wurzeln ist
\begin{equation}
        D>0,
        \qquad a<0,
        \qquad a^2-4c>0.
\tag{6}
\end{equation}
In allen anderen F\"allen, wo $D$ positiv ist, hat die Gleichung vier imagin\"are Wurzeln.

Wir k\"onnen auch f\"ur $D=0$, also im Falle der Gleichheit zweier Wurzeln, die Discussion vollst\"andig durchf\"uhren.

Nehmen wir $\gamma=\delta$ an, so folgt aus (2), da wir die Annahme $\alpha+\beta+\gamma+\delta=0$ gemacht haben,
\[
        2v=2w=\alpha-\beta,
        \qquad
        2u=\alpha+\beta-2\gamma=2(\alpha+\beta),
\]
und daraus nach (3)
\begin{equation}
\begin{aligned}
        -4a&=2(\alpha+\beta)^2+(\alpha-\beta)^2,\cr
        16(a^2-4c)&=(\alpha-\beta)^2[8(\alpha+\beta)^2+(\alpha-\beta)^2].
\end{aligned}
\tag{7}
\end{equation}
Um zun\"achst den Fall zu erledigen, dass auch $\alpha=\beta$ ist, so erhalten wir aus (7) daf\"ur die Bedingung $a^2-4c=0$, und die erste Gleichung (7) zeigt, dass $a$ negativ ist, wenn $\alpha$ und $\gamma$ reell, dagegen positiv, wenn $\alpha$ und $\gamma$ conjugirt (und dann wegen $\alpha+\gamma=0$ rein imagin\"ar) sind. Daraus folgt:

Die biquadratische Gleichung hat zwei Paare gleicher, reeller Wurzeln, wenn
\begin{equation}
        D=0,
        \qquad a<0,
        \qquad a^2-4c=0
\tag{8}
\end{equation}
und zwei Paare gleicher imagin\"arer Wurzeln, wenn
\begin{equation}
        D=0,
        \qquad a>0,
        \qquad a^2-4c=0.
\tag{9}
\end{equation}

Ist aber $\alpha$ von $\beta$ verschieden, so muss $\gamma$ reell sein und (7) zeigt, dass, wenn $\alpha$ und $\beta$ reell sind, $a$ negativ und $a^2-4c$ positiv ist, dass dagegen, wenn $\alpha$ und $\beta$ conjugirt imagin\"ar sind, entweder $a$ positiv oder $a^2-4c$ negativ sein muss, da in diesem Falle $(\alpha-\beta)^2$ negativ ist, und wenn $2(\alpha+\beta)^2+(\alpha-\beta)^2$ positiv ist, jedenfalls auch $8(\alpha+\beta)^2+(\alpha-\beta)^2$ positiv sein muss.

Wir k\"onnen also (6) dahin erg\"anzen, dass wir sagen:

Die nothwendige und hinreichende Bedingung f\"ur die Existenz von vier reellen Wurzeln, von denen auch zwei (aber nicht mehr) einander gleich sein k\"onnen, ist
\begin{equation}
        D\geq 0,
        \qquad a<0,
        \qquad a^2-4c>0.
\tag{10}
\end{equation}
Die nothwendige und hinreichende Bedingung f\"ur drei gleiche Wurzeln, die nothwendig alle reell sind, ist nach (2)
\[
        u^2=v^2=w^2,
\]
es muss also die cubische Resolvente (3) drei gleiche Wurzeln haben, und daf\"ur sind nach \S\,77 (1) die nothwendigen und hinreichenden Bedingungen
\begin{equation}
        a^2+12c=0,
        \qquad
        2a^3-8ac+9b^2=0.
\tag{11}
\end{equation}
Nach \S\,47, (15), (16) sind die beiden Invarianten $A$ und $B$ der biquadratischen Gleichung
\[
        A=a^2+12c,
        \qquad
        B=2a^3-72ac+27b^2.
\]
Also k\"onnen wir die Gleichungen (11) auch so schreiben
\[
        A=0,
        \qquad B+4aA=0,
\]
so dass (11) gleichbedeutend ist mit
\[
        A=0,
        \qquad B=0.
\]
Durch Elimination von $c$ kann man aus (11) auch noch die Gleichung ableiten:
\begin{equation}
        8a^3+27b^2=0.
\tag{12}
\end{equation}
Endlich ist noch die Bedingung f\"ur die Gleichheit aller vier Wurzeln
\begin{equation}
        a=0,
        \qquad b=0,
        \qquad c=0.
\tag{13}
\end{equation}

Man kann diese Verh\"altnisse sehr anschaulich machen durch eine geometrische Deutung, und wenn auch geometrische Betrachtungen nicht eigentlich in unserem Plane liegen, so wollen wir doch nicht unterlassen, den Leser darauf gelegentlich hinzuweisen.

Deutet man $a,b,c$ als rechtwinklige Coordinaten im Raume, so ist jeder Raumpunkt als Tr\"ager einer gewissen biquadratischen Gleichung von der Form (1), $f(x)=0$, zu betrachten; alle Gleichungen, die eine bestimmte Zahl $x$ zur Wurzel haben, werden durch Punkte einer Ebene $[f(x)=0]$ repr\"asentirt. Die Gleichung $D=0$ ist die Gleichung einer krummen Oberfl\"ache (f\"unften Grades), der Discriminantenfl\"ache, die von den Schnittlinien der Ebenen $f(x)=0$, $f'(x)=0$ erzeugt wird, und also eine abwickelbare Fl\"ache ist. Sie ist die Einh\"ullende aller Ebenen $f(x)=0$.

Die Fl\"ache hat eine aus zwei Zweigen bestehende R\"uckkehrkante, die durch die Gleichungen (11), (12) dargestellt ist, und eine Doppellinie, die durch die Gleichungen $b=0$, $a^2-4c=0$ bestimmt ist und also die Gestalt einer Parabel hat. Auf der Seite der negativen $a$ durchsetzt sich in dieser Parabel die Fl\"ache selbst. Auf der Seite der positiven $a$ setzt sich die Parabel als isolirte Linie fort.

Die Discriminantenfl\"ache theilt den ganzen Raum in drei F\"acher, von denen zwei nur l\"angs der Doppelparabel zusammenh\"angen, und diese F\"acher enthalten die Punkte, denen keine, zwei und vier reelle Wurzeln entsprechen. Nennen wir f\"ur den Augenblick diese drei F\"acher $[0]$, $[2]$, $[4]$, so grenzt $[0]$ an $[4]$ nur l\"angs der Doppelparabel, w\"ahrend $[0]$ sowohl als $[4]$ l\"angs der Fl\"achentheile an $[2]$ grenzen. Wir wollen mit $[0,2]$ und $[4,2]$ die Grenzfl\"achen von $[0]$, $[2]$ und von $[0]$, $[4]$ bezeichnen. Der isolirte Theil der Parabel setzt sich in das Innere von $[0]$ fort. Die R\"uckkehrkanten liegen auf dem Theil der Fl\"ache, der $[4]$ von $[2]$ scheidet.

Im Coordinatenanfang stossen die drei Raumtheile und ihre Grenzfl\"achen und Grenzlinien zusammen. Die Punkte der Fl\"achentheile $[0,2]$ stellen Gleichungen mit zwei gleichen und zwei imagin\"aren Wurzeln dar, die Punkte von $[2,4]$ Gleichungen mit zwei gleichen und zwei davon verschiedenen reellen Wurzeln.

Auf der Doppelparabel finden zweimal zwei gleiche Wurzeln statt, und zwar auf dem Theil, in dem $[0]$ an $[4]$ grenzt, reelle, in dem isolirten Theil imagin\"are.


\begin{center}
\begin{tikzpicture}[scale=0.95, line cap=round, line join=round]
  \fill[pattern=north east lines, pattern color=black!65]
    (-3.2,-1.15) .. controls (-3.45,-0.2) and (-3.35,0.7) .. (-2.55,1.05)
    .. controls (-1.95,1.25) and (-1.45,0.95) .. (-1.25,0.25)
    -- (-1.05,-1.45) -- (-2.8,-1.45) -- cycle;
  \draw (-3.2,-1.15) .. controls (-3.45,-0.2) and (-3.35,0.7) .. (-2.55,1.05)
        .. controls (-1.95,1.25) and (-1.45,0.95) .. (-1.25,0.25);
  \draw (-3.2,-1.15) -- (-2.8,-1.45) -- (-1.05,-1.45) -- (-1.25,0.25);
  \draw (-1.25,0.25) .. controls (-0.25,0.95) and (0.55,1.95) .. (1.05,3.15);
  \draw (-1.25,0.25) .. controls (-0.10,0.25) and (0.70,-0.05) .. (1.95,-1.15);
  \draw (1.05,3.15) .. controls (1.65,1.55) and (2.25,0.95) .. (3.25,3.05);
  \draw (1.95,-1.15) .. controls (2.40,0.40) and (2.85,1.65) .. (3.25,3.05);
  \draw (3.25,3.05) -- (3.45,-1.55) -- (1.95,-1.15);
  \draw[dashed] (1.05,3.15) -- (1.10,-1.10);
  \draw[dashed] (3.25,3.05) -- (1.95,-1.15);
  \draw[dashed] (-1.25,0.25) .. controls (0.40,0.30) and (1.90,0.05) .. (2.55,0.75)
        .. controls (2.82,1.05) and (2.95,1.55) .. (3.25,3.05);
  \draw[dashed] (-1.05,-1.45) .. controls (0.05,-1.35) and (1.05,-1.25) .. (1.95,-1.15);
  \draw (1.95,-1.15) -- (3.45,-1.55);
\end{tikzpicture}
\\[0.4em]Fig. 5.
\end{center}


Die Punkte der R\"uckkehrkanten repr\"asentiren Gleichungen mit drei gleichen Wurzeln und der Coordinaten-Anfangspunkt die Gleichung mit vier gleichen Wurzeln.

Auf die Discriminantenfl\"ache und ihre Bedeutung f\"ur die Discussion der biquadratischen Gleichung hat zuerst Kronecker hingewiesen. Ein Modell der Fl\"ache ist von Kerschensteiner construirt.

Der analytische Ausdruck f\"ur die Bedingungen der Realit\"at der Wurzeln l\"asst sich in eine Gestalt bringen, in der nur die Covarianten der biquadratischen Form vorkommen.

Den Ausgangspunkt dazu bilden die Ausdr\"ucke f\"ur die Functionen $\psi_1,\psi_2,\psi_3$, die \S\,66, (6) gegeben sind. Wenn wir in jenen Ausdr\"ucken $\alpha$ mit $\beta$ vertauschen, so gehen $\psi_1,\psi_2,\psi_3$ in $\psi_1,-\psi_3,-\psi_2$ \"uber, und wenn wir gleichzeitig $\alpha$ mit $\beta$ und $\gamma$ mit $\delta$ vertauschen, $\psi_1,\psi_2,\psi_3$ in $\psi_1,-\psi_2,-\psi_3$. Daraus ergiebt sich Folgendes:

Setzt man f\"ur die Variable $x$ einen beliebigen reellen Werth und sind die Wurzeln $\alpha,\beta,\gamma,\delta$ der biquadratischen Form alle reell, so sind auch $\psi_1,\psi_2,\psi_3$ reell.

Sind $\alpha,\beta$ conjugirt imagin\"ar, $\gamma,\delta$ reell, so ist die Vertauschung von $\alpha$ und $\beta$ gleichbedeutend mit der Vertauschung von $\ii$ mit $-\ii$. Es sind also in diesem Falle $\psi_1$ reell, $\psi_2,\psi_3$ conjugirt imagin\"ar.

Sind endlich $\alpha,\beta$ und $\gamma,\delta$ zwei Paare conjugirt imagin\"arer Wurzeln, so werden $\alpha$ mit $\beta$ und $\gamma$ mit $\delta$ vertauscht, wenn $\ii$ in $-\ii$ \"ubergeht; also sind in diesem Falle $\psi_1$ reell, $\psi_2,\psi_3$ rein imagin\"ar.

Sind vier Wurzeln reell, so werden $\psi_1^2,\psi_2^2,\psi_3^2$ reell und positiv. Sind zwei Wurzeln reell, so ist $\psi_1^2$ reell und positiv, $\psi_2^2,\psi_3^2$ sind conjugirt imagin\"ar. Sind vier Wurzeln imagin\"ar, so ist $\psi_1^2$ positiv, $\psi_2^2,\psi_3^2$ sind reell und negativ.

Nun haben wir im \S\,66, (7) die cubische Gleichung aufgestellt, deren Wurzeln $\psi_1^2,\psi_2^2,\psi_3^2$ sind, und haben auch ihre Discriminante gebildet. Im vorigen Paragraphen haben wir ein Kennzeichen f\"ur die Anzahl der positiven Wurzeln einer cubischen Gleichung kennen gelernt, woraus sich folgendes Resultat ergiebt:

Ist $D<0$, so hat die biquadratische Gleichung zwei reelle und zwei imagin\"are Wurzeln.

Ist $D>0$, so m\"ussen, wenn vier reelle Wurzeln vorhanden sein sollen, in der Reihe
\[
        1,
        \qquad H,
        \qquad H^2-16Af^2,
        \qquad -I^2
\]
drei Zeichenwechsel vorkommen, d. h. es muss
\[
        H<0,
        \qquad
        H^2-16Af^2>0
\]
sein. Bei den drei anderen noch m\"oglichen Zeichencombinationen sind alle vier Wurzeln imagin\"ar. Hierbei aber kann f\"ur $x$ ein beliebiger reeller Werth gesetzt werden.

\sect{79}{Die Bezoutiante und ihre Bedeutung f\"ur die Wurzelrealit\"at}

F\"ur die Untersuchung der Realit\"at der Wurzeln einer beliebigen reellen Gleichung in allgemeineren F\"allen kann die Function mit Nutzen angewandt werden, die wir im \S\,73 als Bezoutiante der Gleichung bezeichnet haben.

Ist
\begin{equation}
        f(x)=x^n+a_1x^{n-1}+a_2x^{n-2}+\cdots+a_n=0
\tag{1}
\end{equation}
irgend eine reelle Gleichung $n$ten Grades, so war die Bezoutiante folgendermaassen definirt. Man setze
\begin{equation}
        y=t_{n-1}f_0(x)+t_{n-2}f_1(x)+\cdots+t_0f_{n-1}(x),
\tag{2}
\end{equation}
worin die $t_0,t_1,\ldots,t_{n-1}$ unbestimmte Variable bedeuten und
\[
 f_0(x)=1,
 \qquad f_1(x)=x+a_1,
 \qquad f_2(x)=x^2+a_1x+a_2,
 \quad\ldots
\]
Es ist dann [\S\,73, (10)]
\begin{equation}
        S(y^2)={1\over n}[S(y)]^2+B,
\tag{3}
\end{equation}
wo $B$ eine quadratische Form der $n-1$ Variablen $t_0,t_1,\ldots,t_{n-2}$ ist mit Coefficienten, die rational aus den Coefficienten $a_i$ zusammengesetzt sind. Diese Function $B$ haben wir als Bezoutiante definirt und f\"ur die F\"alle $n=4$ und $n=5$ wirklich gebildet.

Nun sind in der Summe auf der linken Seite von (3)
\begin{equation}
        y_1^2+y_2^2+\cdots+y_n^2
\tag{4}
\end{equation}
die $y_1,y_2,\ldots,y_n$ lineare Functionen der Variablen $t_0,t_1,\ldots,t_{n-1}$, und wir haben also in der Formel (3) eine Transformation der quadratischen Form
\begin{equation}
        {1\over n}[S(y)]^2+B=\Phi(t_0,t_1,\ldots,t_{n-1})
\tag{5}
\end{equation}
auf eine Summe von $n$ Quadraten.

Wir machen f\"urs erste keinerlei beschr\"ankende Voraussetzungen \"uber Gleichheit oder Verschiedenheit der Wurzeln von $f(x)$ und m\"ussen nun zun\"achst untersuchen, inwieweit die linearen Functionen $y_1,y_2,\ldots,y_n$ von einander unabh\"angig sind. In dieser Beziehung gilt der Satz:

Sind $x_1,x_2,\ldots,x_\nu$ von einander verschiedene Wurzeln der Gleichung (1), so sind $y_1,y_2,\ldots,y_\nu$ linear unabh\"angig.

Denn angenommen, es existire zwischen $y_1,y_2,\ldots,y_\nu$ eine in Bezug auf $t$ identische lineare Relation
\[
        \alpha_1y_1+\alpha_2y_2+\cdots+\alpha_\nu y_\nu=0,
\]
deren Coefficienten $\alpha$ von den $t$ unabh\"angig und nicht alle gleich Null sind, so w\"urde diese in die folgenden Gleichungen zerfallen
\[
\begin{array}{rcl}
\alpha_1f_0(x_1)+\alpha_2f_0(x_2)+\cdots+\alpha_\nu f_0(x_\nu)&=&0,\cr
\alpha_1f_1(x_1)+\alpha_2f_1(x_2)+\cdots+\alpha_\nu f_1(x_\nu)&=&0,\cr
\multicolumn{3}{c}{\cdots}\cr
\alpha_1f_{n-1}(x_1)+\alpha_2f_{n-1}(x_2)+\cdots+\alpha_\nu f_{n-1}(x_\nu)&=&0.
\end{array}
\]
Da nun $\nu\leq n$ ist, so muss die Determinante der $\nu$ ersten von diesen Gleichungen verschwinden, also
\[
\begin{vmatrix}
 f_0(x_1)&f_1(x_1)&\cdots&f_{\nu-1}(x_1)\cr
 f_0(x_2)&f_1(x_2)&\cdots&f_{\nu-1}(x_2)\cr
 \vdots&\vdots&&\vdots\cr
 f_0(x_\nu)&f_1(x_\nu)&\cdots&f_{\nu-1}(x_\nu)
\end{vmatrix}=0,
\]
oder
\[
\begin{vmatrix}
 1&x_1+a_1&x_1^2+a_1x_1+a_2&\cdots\cr
 1&x_2+a_1&x_2^2+a_1x_2+a_2&\cdots\cr
 \vdots&\vdots&\vdots&\cr
 1&x_\nu+a_1&x_\nu^2+a_1x_\nu+a_2&\cdots
\end{vmatrix}=0.
\]
Diese Determinante l\"asst sich aber durch wiederholte Anwendung des Satzes von der Addition der Colonnen \S\,22, (VII) auf die Form bringen
\[
\begin{vmatrix}
1&x_1&x_1^2&\cdots&x_1^{\nu-1}\cr
1&x_2&x_2^2&\cdots&x_2^{\nu-1}\cr
\vdots&\vdots&\vdots&&\vdots\cr
1&x_\nu&x_\nu^2&\cdots&x_\nu^{\nu-1}
\end{vmatrix}
=\prod_{1\leq r<s\leq \nu}(x_r-x_s),
\]
und kann also nicht verschwinden, wenn, wie vorausgesetzt war, die $x_1,x_2,\ldots,x_\nu$ verschieden sind.

Wenn wir daher in der Summe (4) alle unter einander gleichen Glieder zusammenfassen, so bleiben so viele Quadrate linear unabh\"angiger Functionen \"ubrig, als die Gleichung (1) verschiedene Wurzeln hat.

Wenn nun $x_1$ eine reelle Wurzel von (1) ist, so ist $y_1^2$ ein positives Quadrat in der Summe (4). Sind aber $x_1$ und $x_2$ conjugirt imagin\"ar, so zerfallen auch $y_1$ und $y_2$ in zwei conjugirt imagin\"are Bestandtheile
\[
        y_1=u+v\ii,
        \qquad
        y_2=u-v\ii,
\]
wo $u$ und $v$ lineare reelle Functionen von den $t$ sind. Demnach wird
\[
        y_1^2+y_2^2=2u^2-2v^2.
\]

Es ist also durch (4) die durch (5) definirte Function $\Phi(t)$ in eine Summe von Quadraten zerlegt, deren Anzahl gleich der Zahl der verschiedenen Wurzeln von $f(x)=0$ ist, und unter denen so viele negative sind, als unter diesen von einander verschiedenen Wurzeln Paare imagin\"arer Wurzeln vorkommen.

Diese Anzahlen bleiben aber nach dem Tr\"agheitsgesetz der quadratischen Formen (\S\,58) bei jeder anderen linearen reellen Transformation der quadratischen Form in eine Summe von Quadraten dieselben. Nun ist $S(y)$ eine reelle lineare Function der $t$, und $B$ enth\"alt die Variable $t_{n-1}$ nicht mehr; hiernach ergiebt sich aus der Formel (5) der Satz:

I. Wenn die Bezoutiante $B$ durch reelle lineare Transformation in eine Summe von $\pi$ positiven und $\nu$ negativen Quadraten, die nicht auf eine kleinere Zahl reducirt werden k\"onnen, zerlegt ist, so ist $\pi+\nu+1$ die Zahl der verschiedenen Wurzeln von $f(x)=0$; und $\nu$ ist die Anzahl der darunter enthaltenen Paare conjugirt imagin\"arer Wurzeln, $\pi-\nu+1$ die der reellen.

Ist die Determinante von $B$, also auch die Discriminante von $f(x)$ (\S\,73) von Null verschieden, so ist $\pi+\nu+1=n$ und $\nu$ ist kleiner oder h\"ochstens gleich $\frac12n$.

Die Bezoutiante ist also eine quadratische Form von einer besonderen Natur, die sich eben darin ausspricht, dass unter den Quadraten, in die sie sich zerlegen l\"asst, h\"ochstens $\frac12n$ negative vorkommen k\"onnen.

\sect{80}{Die Tr\"agheit der Formen zweiten Grades}

Durch den Satz des vorigen Paragraphen sind wir auf die Untersuchung der homogenen Functionen zweiten Grades mit reellen Coefficienten hingewiesen. Ueber diese Functionen haben wir in \S\,58 unter dem Namen des Tr\"agheitsgesetzes einen Satz kennen gelernt, der f\"ur das Folgende die Grundlage bildet. Der Satz bestand darin, dass, wie man auch eine quadratische Form von $m$ Ver\"anderlichen in eine Summe von positiven und negativen Quadraten von linear unabh\"angigen linearen Functionen transformiren mag, was auf unendlich viele verschiedene Arten m\"oglich ist, die Anzahl der positiven und ebenso die der negativen Quadrate immer dieselbe ist.

Bezeichnen wir mit $\pi$ die Anzahl der positiven und mit $\nu$ die Anzahl der negativen Quadrate, so ist $\pi+\nu$ h\"ochstens gleich $m$.

Wenn wir den Unterschied $m-\pi-\nu$ mit $\rho$ bezeichnen, so kann, wenn wir die allgemeine quadratische Form von $m$ Variablen als Summe von $m$ Quadraten darstellen, $\rho$ als die Anzahl der Quadrate mit verschwindenden Coefficienten bezeichnet werden.

Die drei Zahlen $\pi,\nu,\rho$, von denen keine negativ sein kann und deren Summe gleich der Anzahl der Variablen ist,
\begin{equation}
        \pi+\nu+\rho=m,
\tag{1}
\end{equation}
sind also f\"ur eine bestimmte quadratische Form unver\"anderlich.

Wir haben nach Mitteln zu suchen, um aus den Coefficienten der quadratischen Form die Zahlen $\pi,\nu,\rho$ zu ermitteln. Wenden wir diese Mittel auf die Bezoutiante an, so erhalten wir Kennzeichen, um aus den Coefficienten einer Gleichung auf die Anzahl ihrer reellen Wurzeln zu schliessen.

Es sei also jetzt, wie in \S\,56,
\begin{equation}
        \varphi(x_1,x_2,\ldots,x_m)=\sum_{1}^{m}a_{i,k}x_ix_k,
        \qquad a_{i,k}=a_{k,i},
\tag{2}
\end{equation}
eine quadratische Form von $m$ Ver\"anderlichen mit reellen Coefficienten $a_{i,k}$.

Die Determinante
\begin{equation}
        R=\sum \pm a_{1,1}a_{2,2}\cdots a_{m,m}
\tag{3}
\end{equation}
zeigt durch ihr Verschwinden an, dass $\varphi$ durch lineare Transformation in eine Function von weniger als $m$ Variablen transformirt werden kann, dass also $\rho$ gr\"osser als Null ist.

Die Determinante $R$ ist eine symmetrische Determinante. Unter ihren Unterdeterminanten verschiedener Ordnung kommen gewisse vor, die wieder symmetrische Determinanten sind, n\"amlich die, die man aus $R$ erh\"alt, wenn man Zeilen und Colonnen, die sich in Diagonalgliedern schneiden, ausstreicht, die also, wenn $\alpha,\beta,\gamma,\ldots$ irgend welche unter den Indices $1,2,\ldots,m$ bedeuten, durch
\begin{equation}
        \sum\pm a_{\alpha,\alpha}a_{\beta,\beta}a_{\gamma,\gamma}\cdots
\tag{4}
\end{equation}
zu bezeichnen sind. Diese wollen wir die Haupt-Unterdeterminanten von $R$ nennen.

Wir wollen die Anzahl der Haupt-Unterdeterminanten bestimmen.

Die Anzahl der $\mu$-reihigen Haupt-Unterdeterminanten ist gleich der Anzahl der Arten, wie man aus der Reihe der Indices $1,2,3,\ldots,m$ Gruppen von $\mu$ verschiedenen ausw\"ahlen kann, also gleich der Anzahl der Combinationen von $m$ Elementen zu je $\mu$, und diese Zahl ist gleich dem Binomialcoefficienten $B^m_\mu$. Die Anzahl aller Haupt-Unterdeterminanten ist also, wenn wir die Determinante $R$ selbst und ausserdem noch die Einheit als eine nullreihige Determinante mitz\"ahlen,
\[
        1+m+{m(m-1)\over 1\cdot 2}+\cdots+m+1=2^m.
\]

Es ist aber meist nur ein kleiner Theil von diesen wirklich zu ber\"ucksichtigen.

Die Indices $1,2,3,\ldots,m$ lassen sich auf $\Pi(m)=1\cdot2\cdot3\cdots m$ verschiedene Arten anordnen; wenn wir mit irgend einer dieser Anordnungen die Determinante bilden
\[
\begin{vmatrix}
 a_{1,1} & a_{1,2} & a_{1,3} & \cdots & a_{1,m}\cr
 a_{2,1} & a_{2,2} & a_{2,3} & \cdots & a_{2,m}\cr
 a_{3,1} & a_{3,2} & a_{3,3} & \cdots & a_{3,m}\cr
 \vdots  & \vdots  & \vdots  &        & \vdots\cr
 a_{m,1} & a_{m,2} & a_{m,3} & \cdots & a_{m,m}
\end{vmatrix},
\]
so l\"asst sich eine Reihe von Haupt-Unterdeterminanten daraus ableiten, wie es durch die Striche angedeutet ist, d. h. so, dass man jede vorhergehende aus der nachfolgenden erh\"alt, indem man die letzte Zeile und Colonne wegl\"asst; es ist also
\[
        R_0=1,
        \qquad R_1=a_{1,1},
        \qquad R_2=a_{1,1}a_{2,2}-a_{1,2}^2,
\]
\[
        R_3=\begin{vmatrix}
        a_{1,1}&a_{1,2}&a_{1,3}\cr
        a_{2,1}&a_{2,2}&a_{2,3}\cr
        a_{3,1}&a_{3,2}&a_{3,3}
        \end{vmatrix},
        \qquad\ldots,
        \qquad R_m=R.
\]
Ein solches System soll, wenn es in absteigender Reihe
\[
        R_m,R_{m-1},\ldots,R_1,R_0
\]
geordnet ist, eine Kette von Haupt-Unterdeterminanten heissen. Solcher Ketten lassen sich $\Pi(m)$ verschiedene bilden, in denen allen das erste Glied $R$, das letzte Glied $1$ ist.

\sect{81}{Quadratische Formen mit verschwindender Determinante}

Wenn die in (3), \S\,80 definirte Determinante $R$ der quadratischen Form $\varphi(x_1,x_2,\ldots,x_m)$ verschwindet, so l\"asst sich, wie wir schon im \S\,57 gesehen haben, $\varphi$ durch weniger als $m$ von einander unabh\"angige lineare Functionen von $x$ ausdr\"ucken. Wenn wir mit $m-\rho$ die kleinste Zahl von Variablen bezeichnen, durch die sich $\varphi$ ausdr\"ucken l\"asst, so hat $\rho$ dieselbe Bedeutung, wie im vorigen Paragraphen.

Wenn sich nun unter den Haupt-Unterdeterminanten von $R$ eine $k$-reihige findet, die von Null verschieden ist, etwa
\[
        R_k=\sum \pm a_{1,1}a_{2,2}\cdots a_{k,k},
\]
so ist, wenn wir
\[
        x_{k+1}=0,
        \qquad x_{k+2}=0,
        \qquad\ldots\qquad x_m=0
\]
setzen,
\[
        \varphi(x_1,x_2,\ldots,x_k,0,0,0)
\]
eine quadratische Form von $k$ Variablen $x_1,x_2,\ldots,x_k$, deren Determinante $R_k$ von Null verschieden ist, die sich sonach nicht durch weniger als $k$ Variable ausdr\"ucken l\"asst. Um so weniger kann also bei unbeschr\"ankt ver\"anderlichen $x_1,x_2,\ldots,x_m$ die Function $\varphi$ von weniger als $k$ Variablen abh\"angen, und es folgt
\begin{equation}
        \rho\leq m-k.
\tag{1}
\end{equation}

Wir nehmen nun an, die Form $\varphi(x_1,x_2,\ldots,x_m)$ lasse sich durch $k$ von einander unabh\"angige lineare Formen von $x$ ausdr\"ucken, die wir mit
\begin{equation}
\begin{aligned}
 \xi_1&=\beta^{(1)}_1x_1+\beta^{(1)}_2x_2+\cdots+\beta^{(1)}_kx_k+\beta^{(1)}_{k+1}x_{k+1}+\cdots+\beta^{(1)}_mx_m,\cr
 \cdots&\cr
 \xi_k&=\beta^{(k)}_1x_1+\beta^{(k)}_2x_2+\cdots+\beta^{(k)}_kx_k+\beta^{(k)}_{k+1}x_{k+1}+\cdots+\beta^{(k)}_mx_m
\end{aligned}
\tag{2}
\end{equation}
bezeichnen wollen.

Da $\xi_1,\xi_2,\ldots,\xi_k$ von einander linear unabh\"angig sein sollen, so muss unter den aus der Matrix
\[
\begin{array}{cccc}
 \beta^{(1)}_1&\beta^{(1)}_2&\cdots&\beta^{(1)}_m\cr
 \cdots&\cdots&&\cdots\cr
 \beta^{(k)}_1&\beta^{(k)}_2&\cdots&\beta^{(k)}_m
\end{array}
\]
zu bildenden $k$-reihigen Determinanten wenigstens eine von Null verschieden sein. Denn w\"aren sie alle gleich Null, so liessen sich die Unbekannten $h_1,h_2,\ldots,h_k$, die nicht alle verschwinden sollen, aus den Gleichungen
\[
        h_1\beta^{(1)}_s+h_2\beta^{(2)}_s+\cdots+h_k\beta^{(k)}_s=0,
        \qquad s=1,2,3,\ldots,m,
\]
bestimmen (\S\,24), und die $\xi_1,\xi_2,\ldots,\xi_k$ w\"urden einer Gleichung
\[
        h_1\xi_1+h_2\xi_2+\cdots+h_k\xi_k=0
\]
gen\"ugen, also nicht unabh\"angig sein.

Wir k\"onnen aber, ohne die Allgemeinheit zu beschr\"anken, annehmen, dass die Determinante
\[
        \sum\pm \beta^{(1)}_1\beta^{(2)}_2\cdots\beta^{(k)}_k
\]
von Null verschieden sei, und dann k\"onnen wir die Gleichungen (2) in Bezug auf die Variablen $x_1,x_2,\ldots,x_k$ aufl\"osen.

Diesen Aufl\"osungen geben wir die Form
\begin{equation}
\begin{aligned}
 x_1&=y_1+\alpha^{(1)}_{k+1}x_{k+1}+\cdots+\alpha^{(1)}_mx_m,\cr
 x_2&=y_2+\alpha^{(2)}_{k+1}x_{k+1}+\cdots+\alpha^{(2)}_mx_m,\cr
 \cdots&\cr
 x_k&=y_k+\alpha^{(k)}_{k+1}x_{k+1}+\cdots+\alpha^{(k)}_mx_m,
\end{aligned}
\tag{3}
\end{equation}
worin die $y_s$ lineare homogene Verbindungen der $\xi_s$ sind, die aus den Gleichungen
\[
\begin{array}{rcl}
 \beta^{(1)}_1y_1+\beta^{(1)}_2y_2+\cdots+\beta^{(1)}_ky_k&=&\xi_1,\cr
 \beta^{(2)}_1y_1+\beta^{(2)}_2y_2+\cdots+\beta^{(2)}_ky_k&=&\xi_2,\cr
 \multicolumn{3}{c}{\cdots}\cr
 \beta^{(k)}_1y_1+\beta^{(k)}_2y_2+\cdots+\beta^{(k)}_ky_k&=&\xi_k
\end{array}
\]
bestimmt werden, w\"ahrend die $\alpha$ neue Coefficienten sind, die in einer leicht zu \"ubersehenden Weise aus den $\beta$ abgeleitet werden, auf deren Bildung es uns hier nicht weiter ankommt.

Nun l\"asst sich nach der Voraussetzung die Function $\varphi(x_1,x_2,\ldots,x_m)$ durch $\xi_1,\xi_2,\ldots,\xi_k$, also auch durch $y_1,y_2,\ldots,y_k$ allein ausdr\"ucken. Bezeichnen wir diesen Ausdruck mit $\Phi(y_1,y_2,\ldots,y_k)$, so haben wir durch Vermittelung von (3) die Identit\"at
\begin{equation}
        \varphi(x_1,x_2,\ldots,x_k,x_{k+1},\ldots,x_m)=\Phi(y_1,y_2,\ldots,y_k),
\tag{4}
\end{equation}
und wenn wir darin $x_{k+1},x_{k+2},\ldots,x_m=0$ setzen,
\begin{equation}
        \Phi(x_1,x_2,
        \ldots,x_k)=\varphi(x_1,x_2,\ldots,x_k,0,0,\ldots,0),
\tag{5}
\end{equation}
wodurch, da es auf die Bezeichnung der Variablen nicht ankommt, $\Phi$ vollst\"andig bestimmt ist; wir k\"onnen daher setzen
\begin{equation}
\begin{aligned}
        \Phi(y_1,y_2,\ldots,y_k)&=\varphi(y_1,y_2,\ldots,y_k,0,0,\ldots,0)\cr
        &=\varphi(x_1,x_2,\ldots,x_m)',
\end{aligned}
\tag{6}
\end{equation}
$\Phi$ entsteht also aus $\varphi$ dadurch, dass man die $m-k$ letzten Variablen (bei der hier gew\"ahlten Anordnung) gleich Null setzt:
\begin{equation}
        \Phi(x_1,\ldots,x_k)=\sum_{1}^{k}a_{r,s}x_rx_s.
\tag{7}
\end{equation}

Die Determinante von $\Phi$ ist eine der Haupt-Unterdeterminanten von $R$, und zwar erh\"alt man sie, indem man in $R$ die $m-k$ letzten Zeilen und Colonnen wegstreicht.

Nehmen wir an, dass $\varphi$ sich nicht durch noch weniger als $k$ Variable ausdr\"ucken l\"asst, dass also $k=m-\rho$ sei, so kann die Determinante dieser Function $\Phi$ nicht verschwinden.

Hieraus und aus dem zuvor Bewiesenen ergiebt sich nun der Satz:

II. Wenn alle Haupt-Unterdeterminanten von $R$ von mehr als $k$ Reihen verschwinden, w\"ahrend unter den Haupt-Unterdeterminanten von $k$ Reihen wenigstens eine von Null verschieden ist, so ist
\[
        \rho=m-k.
\]
Denn aus (7) folgt, dass, wenn $\rho=m-k$ ist, wenigstens eine $k$-reihige Haupt-Unterdeterminante von Null verschieden sein muss, und aus (1), dass, wenn auch noch eine Haupt-Unterdeterminante von mehr als $k$ Reihen von Null verschieden ist,
\[
        \rho<m-k
\]
ist.

Sind $\pi,\nu,\rho$ die Anzahl der positiven, negativen, verschwindenden Quadrate, in die sich $\varphi$ zerlegen l\"asst, so sind $\pi,\nu$ die Anzahlen der positiven und negativen Quadrate, in die sich die durch (7) bestimmte Form $\Phi$ zerlegen l\"asst, und die Bestimmung der Zahlen $\pi,\nu$ braucht also nur noch f\"ur letztere Function, deren Determinante von Null verschieden ist, durchgef\"uhrt zu werden.

\end{document}
